The Aho-Corasick algorithm is a cornerstone of signature-based network
intrusion detection systems and malware analysis tools, because of its linear-time
complexity in the size of the input. Its state-transition dependency, however,
ties the canonical implementation to sequential execution, leaving modern
multi-core processors underused. This work designs, implements, and evaluates
data-parallel scanning strategies built on POSIX Threads for shared-memory
multi-core CPUs. The input text is partitioned into overlapping segments that
worker threads scan concurrently over a shared, strictly read-only automaton,
keeping the hot path free of locks; a flattened, contiguous output layout is
also proposed to eliminate the pointer-chasing cost of emitting matches. The
study deliberately targets the memory-bound regime where the automaton greatly
exceeds the last-level cache, using realistic Snort and Emerging Threats rule
sets (up to roughly $507$~MiB of automaton) over multi-gigabyte corpora on a
hybrid Intel Core i5-1235U processor. The results show that text-level data
parallelism outperforms pattern-set partitioning: search speedups reach up
to $4.79\times$ on twelve threads for a moderate dictionary and $4.52\times$ at
the multi-gigabyte scale, while the flattened layout adds a regime-dependent gain
that grows with memory pressure, up to $+13.0\%$ on a single thread. When the
automaton is forty-two times larger than the cache, however, throughput drops by
$74.0\%$ and the parallel speedup peaks near $2.85\times$ at six threads on the
canonical corpus before regressing, reaching $2.79\times$ at twelve threads on the
multi-gigabyte corpus, exposing memory bandwidth as the physical ceiling. End-to-end accelerations reach up to
$4.48\times$. The central finding is that the dominant barrier to scaling
Aho-Corasick for intrusion detection is microarchitectural, namely cache
overflow and DRAM bandwidth, rather than the matching logic itself.
